% ============================================================
% 4_D_student_ui.tex — 学生端页面详细设计 + 控件 + 校验
% 编写人：D  日期：2026-07-15
% ============================================================

\section{学生端页面详细设计概述}

学生端共涉及8个核心页面/对话框，本章逐一描述每个页面的布局结构、控件清单、数据绑定、交互逻辑及输入校验规则。所有页面遵循Element Plus设计规范，统一使用 \#409EFF（主题蓝）作为主色调，\#67C23A（成功绿）、\#E6A23C（警告橙）、\#F56C6C（危险红）作为辅助功能色。

% ============================================================
\subsection{页面1：登录页（Login.vue）}

\subsubsection{页面布局}
居中双栏卡片布局，整体宽度960px，垂直水平居中。左侧（占40\%宽度）为品牌展示区，显示系统Logo、名称"论文生成系统"、副标题。右侧（占60\%宽度）为登录表单区。

\subsubsection{控件清单}
\begin{table}[htbp]
    \centering
    \caption{登录页控件明细}
    \begin{tabular}{|c|c|c|c|}
        \hline
        \textbf{控件类型} & \textbf{字段名} & \textbf{控件说明} & \textbf{校验规则} \\
        \hline
        el-input & username & 学号输入框 & 必填，纯数字，8-12位 \\
        \hline
        el-input(type=password) & password & 密码输入框，带显示/隐藏图标 & 必填，≥6位 \\
        \hline
        el-checkbox & remember & "记住我"复选框 & 无校验 \\
        \hline
        el-button(type=primary) & — & "登 录"按钮，宽度100\% & 表单校验通过后可用 \\
        \hline
        el-link & — & "还没有账号？立即注册"链接 & — \\
        \hline
    \end{tabular}
\end{table}

\subsubsection{交互逻辑}
\begin{enumerate}
    \item 用户输入学号和密码，点击"登录"按钮或按Enter键提交；
    \item 前端调用表单校验，校验通过后调用 authStore.login(username, password)；
    \item 登录成功后跳转至论文列表页 /papers；登录失败在输入框下方显示红色错误提示；
    \item 若勾选"记住我"，Token存储至localStorage；否则使用sessionStorage。
\end{enumerate}

\subsubsection{输入校验规则}
\begin{lstlisting}[language=JavaScript]
const loginRules = {
    username: [
        { required: true, message: '请输入学号', trigger: 'blur' },
        { pattern: /^\d{8,12}$/, message: '学号为8-12位数字', trigger: 'blur' }
    ],
    password: [
        { required: true, message: '请输入密码', trigger: 'blur' },
        { min: 6, message: '密码长度不少于6位', trigger: 'blur' }
    ]
}
\end{lstlisting}

% ============================================================
\subsection{页面2：注册页（Register.vue）}

\subsubsection{页面布局}
布局风格与登录页一致，右侧表单区域扩展为5个字段。注册成功后自动跳转至登录页并弹出绿色成功提示。

\subsubsection{控件清单}
\begin{table}[htbp]
    \centering
    \caption{注册页控件明细}
    \begin{tabular}{|c|c|c|c|}
        \hline
        \textbf{控件类型} & \textbf{字段名} & \textbf{控件说明} & \textbf{校验规则} \\
        \hline
        el-input & studentId & 学号 & 必填，数字，8-12位，后端唯一校验 \\
        \hline
        el-input & name & 姓名 & 必填，2-20个字符 \\
        \hline
        el-input & email & 邮箱 & 必填，合法邮箱格式 \\
        \hline
        el-input(type=password) & password & 密码 & 必填，≥6位，含字母+数字 \\
        \hline
        el-input(type=password) & confirmPassword & 确认密码 & 必填，与password一致 \\
        \hline
        el-button(type=primary) & — & "注 册"按钮 & 所有校验通过后可用 \\
        \hline
    \end{tabular}
\end{table}

\subsubsection{交互逻辑}
\begin{enumerate}
    \item 用户填写全部字段，前端实时校验各项输入；
    \item 学号输入框在失焦时触发唯一性校验；
    \item 全部校验通过后，提交 POST /api/v1/auth/register；
    \item 注册成功跳转登录页；注册失败显示具体错误。
\end{enumerate}

\subsubsection{输入校验规则}
\begin{lstlisting}[language=JavaScript]
const registerRules = {
    studentId: [
        { required: true, message: '请输入学号', trigger: 'blur' },
        { pattern: /^\d{8,12}$/, message: '学号为8-12位数字', trigger: 'blur' }
    ],
    name: [
        { required: true, message: '请输入姓名', trigger: 'blur' },
        { min: 2, max: 20, message: '姓名为2-20个字符', trigger: 'blur' }
    ],
    email: [
        { required: true, message: '请输入邮箱', trigger: 'blur' },
        { type: 'email', message: '邮箱格式不正确', trigger: 'blur' }
    ],
    password: [
        { required: true, message: '请输入密码', trigger: 'blur' },
        { min: 6, message: '密码长度不少于6位', trigger: 'blur' },
        { pattern: /^(?=.*[a-zA-Z])(?=.*\d)/, message: '密码需包含字母和数字', trigger: 'blur' }
    ],
    confirmPassword: [
        { required: true, message: '请确认密码', trigger: 'blur' },
        { validator: (rule, value, callback) => {
            if (value !== form.password) callback(new Error('两次密码不一致'))
            else callback()
        }, trigger: 'blur' }
    ]
}
\end{lstlisting}

% ============================================================
\subsection{页面3：论文列表页（PaperList.vue）}

\subsubsection{页面布局}
顶部：el-input搜索框 + el-select排序下拉框 + el-button(type=primary)"新建论文"按钮。主体区域：论文卡片网格布局（响应式：≥1200px四列，≥768px两列，<768px单列）。

\subsubsection{控件清单}
\begin{table}[htbp]
    \centering
    \caption{论文列表页控件明细}
    \begin{tabular}{|c|c|c|}
        \hline
        \textbf{控件类型} & \textbf{控件说明} & \textbf{数据绑定} \\
        \hline
        el-input(搜索) & 按论文标题搜索，输入防抖300ms & v-model="searchKeyword" \\
        \hline
        el-select & 排序方式选择 & v-model="sortBy" \\
        \hline
        el-button(type=primary) & 新建论文按钮 & @click前往创建页 \\
        \hline
        PaperCard(自定义) & 论文卡片组件 & v-for循环渲染 \\
        \hline
        el-pagination & 分页器（每页12条） & v-model:current-page \\
        \hline
        el-empty & 空状态插画 & v-if空列表 \\
        \hline
    \end{tabular}
\end{table}

\subsubsection{论文卡片子组件设计}
每张论文卡片（PaperCard）包含：
\begin{itemize}
    \item 论文标题（el-text，最多显示2行，超出省略号截断）；
    \item 模板名称标签（el-tag，type=info，size=small）；
    \item 最后修改时间（相对时间格式："3小时前"、"昨天 14:30"）；
    \item 字数统计（"约12,500字"）；
    \item 操作按钮组：编辑（type=primary, link）、预览（type=success, link）、删除（type=danger, link）。
\end{itemize}

\subsubsection{交互逻辑}
\begin{enumerate}
    \item 页面挂载时调用 GET /api/v1/papers 获取论文列表；
    \item 搜索关键词变更（防抖300ms）或排序方式变更时，重新请求后端接口；
    \item 点击"编辑"按钮，跳转至 /editor/:paperId；
    \item 点击"预览"按钮，在新标签页打开 /preview/:paperId；
    \item 点击"删除"按钮，弹出ElMessageBox二次确认对话框，确认后删除；
    \item 论文卡片hover时浮现浅色阴影和上移动画。
\end{enumerate}

% ============================================================
\subsection{页面4：论文创建向导（PaperCreate.vue）}

\subsubsection{页面布局}
使用el-steps步骤条引导三步创建流程：基本信息 → 选择模板 → 确认创建。

\subsubsection{控件清单}
\begin{table}[htbp]
    \centering
    \caption{论文创建向导控件明细}
    \begin{tabular}{|c|c|c|c|}
        \hline
        \textbf{步骤} & \textbf{控件类型} & \textbf{控件说明} & \textbf{校验规则} \\
        \hline
        第一步 & el-input & 论文标题 & 必填，1-200字符 \\
        \cline{2-4}
        & el-select & 所属学院选择 & 必选 \\
        \hline
        第二步 & TemplateCard(自定义) & 模板卡片列表，单选 & 必选 \\
        \hline
        第三步 & el-descriptions & 汇总展示标题、学院、模板 & — \\
        \cline{2-4}
        & el-button(type=primary) & "确认创建"按钮 & — \\
        \hline
    \end{tabular}
\end{table}

\subsubsection{交互逻辑}
\begin{enumerate}
    \item 步骤条已完成步骤可点击回退，未完成步骤不可跳过；
    \item 第一步选择学院后，第二步自动加载对应模板列表；
    \item 第二步点击模板卡片进行单选，高亮选中状态；
    \item 第三步展示汇总信息，点击"确认创建"提交，成功后自动跳转编辑器。
\end{enumerate}

% ============================================================
\subsection{页面5：论文编辑器主页面（PaperEditor.vue）}

这是学生端最核心、最复杂的页面，采用三栏可调整布局。

\subsubsection{页面布局}
整体高度100vh，三栏布局：
\begin{itemize}
    \item \textbf{顶部工具栏栏（高度48px）}：论文标题显示 + 保存状态指示器 + 预览按钮 + 导出按钮 + 返回列表按钮；
    \item \textbf{左侧章节树面板（默认260px，可拖拽宽度）}：章节树控件 + 添加章节按钮，可收起至40px；
    \item \textbf{中心编辑区（flex:1）}：顶部格式工具栏 + Tiptap富文本编辑区域；
    \item \textbf{右侧辅助面板（默认320px）}：Tab切换"参考文献" / "模板信息"，可收起。
\end{itemize}

\subsubsection{控件清单}

\textbf{顶部工具栏栏：}
\begin{itemize}
    \item el-text(tag="h3") — 论文标题，点击可编辑
    \item SaveIndicator(自定义) — 保存状态："已保存 ✓ 15:30" / "保存中..." / "未保存 ⚠"
    \item el-button(plain) — "预览"按钮，新标签页打开预览
    \item el-button(type=primary) — "导出"按钮，弹出ExportDialog
    \item el-button(text) — "← 返回列表"按钮，含未保存确认
\end{itemize}

\textbf{左侧章节树面板：}
\begin{itemize}
    \item el-tree — 章节树，支持node拖拽排序，自定义节点渲染，右键上下文菜单
    \item el-button(text) — "+ 添加章节"按钮
\end{itemize}

\textbf{章节树右键菜单（自定义el-dropdown）：}
\begin{itemize}
    \item "添加同级章节"、"添加子章节"、"重命名"（章节名称变为el-input）、"删除章节"（含子章节提示）
\end{itemize}

\textbf{中心编辑区格式工具栏：}
\begin{table}[htbp]
    \centering
    \caption{编辑器格式工具栏控件}
    \begin{tabular}{|c|c|c|}
        \hline
        \textbf{分组} & \textbf{控件} & \textbf{Tiptap命令} \\
        \hline
        文本样式 & el-button-group(加粗/斜体/下划线/删除线) & toggleBold/Italic/Underline/Strike \\
        \hline
        标题层级 & el-select(正文/H1/H2/H3/H4) & toggleHeading \\
        \hline
        文字颜色 & el-color-picker(文字色/背景色) & setColor/setHighlight \\
        \hline
        对齐方式 & el-button-group(左/中/右/两端) & setTextAlign \\
        \hline
        列表 & el-button-group(有序/无序) & toggleOrderedList/BulletList \\
        \hline
        缩进 & el-button-group(增/减缩进) & indent/outdent \\
        \hline
        插入 & el-button(引用/图片/表格) & 打开对应面板 \\
        \hline
        撤销/重做 & el-button-group & undo/redo \\
        \hline
    \end{tabular}
\end{table}

\textbf{右侧辅助面板（Tab切换）：}
\begin{itemize}
    \item 参考文献管理Tab：el-table(列表) + el-button(添加) + 行内编辑/删除/引用按钮
    \item 模板信息Tab：当前模板名称 + 样式描述 + 缩略图 + el-button("切换模板")
\end{itemize}

\subsubsection{编辑器核心交互流程}

\begin{enumerate}
    \item \textbf{页面初始化}：获取paperId → 并行请求论文详情和章节树 → 创建Tiptap实例 → 加载首章节 → 注册beforeunload监听
    \item \textbf{章节切换}：点击el-tree节点 → 检查未保存修改 → 执行保存 → 请求新章节内容 → editor.commands.setContent()
    \item \textbf{自动保存（useAutoSave composable）}：监听editor onUpdate → 标记脏数据 → 30秒防抖 → 自动PUT请求 → 更新保存状态
    \item \textbf{章节拖拽}：el-tree draggable → node-drop事件 → PUT提交新顺序 → 失败则回滚
\end{enumerate}

\subsubsection{章节名称校验规则}
\begin{lstlisting}[language=JavaScript]
const sectionNameRules = [
    { required: true, message: '请输入章节名称', trigger: 'blur' },
    { max: 100, message: '章节名称不超过100个字符', trigger: 'blur' },
    { pattern: /^[^\<\>\:\"\/\\\|\?\*]+$/, message: '不能包含特殊字符', trigger: 'blur' }
]
\end{lstlisting}

% ============================================================
\subsection{页面6：参考文献管理面板（ReferencePanel.vue）}

\subsubsection{布局}
侧边面板内垂直布局：顶部"添加文献"按钮，主体el-table（列：序号、作者、标题、年份、操作），底部统计"共N条文献"。

\subsubsection{参考文献表单校验规则}
\begin{lstlisting}[language=JavaScript]
const referenceRules = {
    type: [{ required: true, message: '请选择文献类型', trigger: 'change' }],
    author: [{ required: true, message: '请输入作者', trigger: 'blur' }],
    title: [{ required: true, message: '请输入标题', trigger: 'blur' },
            { max: 500, message: '标题不超过500字符', trigger: 'blur' }],
    journal: [{ required: true, message: '请输入期刊/出版社', trigger: 'blur' }],
    year: [{ required: true, message: '请输入年份', trigger: 'blur' },
           { pattern: /^\d{4}$/, message: '年份为4位数字', trigger: 'blur' }],
    doi: [{ pattern: /^10\.\d+/, message: 'DOI格式不正确', trigger: 'blur' }]
}
\end{lstlisting}

\subsubsection{交互逻辑}
\begin{enumerate}
    \item 面板展开时加载已有文献列表；
    \item 添加/编辑操作弹出el-dialog，表单字段根据文献类型动态显示；
    \item 删除操作弹出确认框，若正文有引用则提示"正文中有N处引用该文献"；
    \item 用户点击文献行"引用"按钮，编辑器光标位置插入引用标记节点。
\end{enumerate}

% ============================================================
\subsection{页面7：导出对话框（ExportDialog.vue）}

\subsubsection{布局}
el-dialog模态对话框（宽度560px），纵向排列导出配置选项。

\subsubsection{控件清单}
\begin{table}[htbp]
    \centering
    \caption{导出对话框控件明细}
    \begin{tabular}{|c|c|c|}
        \hline
        \textbf{控件类型} & \textbf{控件说明} & \textbf{默认值} \\
        \hline
        el-radio-group & 导出格式（DOCX / PDF） & PDF \\
        \hline
        el-checkbox-group & 导出范围（全部/自定义）+ el-tree-select & 全部章节 \\
        \hline
        el-checkbox & "包含封面页" & 勾选 \\
        \hline
        el-checkbox & "包含目录" & 勾选 \\
        \hline
        el-checkbox & "包含参考文献列表" & 勾选 \\
        \hline
        el-progress & 导出进度条（提交后显示） & 0\% \\
        \hline
        el-button(type=primary) & "开始导出" / "下载文件" & 状态切换 \\
        \hline
    \end{tabular}
\end{table}

\subsubsection{交互逻辑}
\begin{enumerate}
    \item 选择格式和范围 → 提交 POST /api/v1/papers/:id/export → 返回taskId；
    \item 前端每2秒轮询 GET /api/v1/export/:taskId/status → 更新进度条；
    \item 完成时按钮变为"下载文件"，触发浏览器下载；
    \item 失败时显示红色错误信息 + "重新导出"按钮。
\end{enumerate}

% ============================================================
\subsection{页面8：个人中心页（Profile.vue）}

\subsubsection{布局}
el-tabs切换三个子面板：个人信息、修改密码、导出历史。el-form表单布局（label-width="100px"）。

\subsubsection{控件摘要}
\begin{itemize}
    \item \textbf{个人信息Tab}：el-input（姓名、邮箱、手机号）+ el-button("保存修改")
    \item \textbf{修改密码Tab}：el-input(type=password)（原密码、新密码、确认新密码）+ el-button("修改密码")
    \item \textbf{导出历史Tab}：el-table（文件名/格式/时间/状态/下载按钮）
\end{itemize}

\subsubsection{校验规则}
\begin{lstlisting}[language=JavaScript]
const profileRules = {
    name: [{ required: true, min: 2, max: 20, message: '姓名为2-20个字符', trigger: 'blur' }],
    email: [{ required: true, type: 'email', message: '邮箱格式不正确', trigger: 'blur' }],
    phone: [{ pattern: /^1[3-9]\d{9}$/, message: '手机号格式不正确', trigger: 'blur' }]
}
const passwordRules = {
    oldPassword: [{ required: true, message: '请输入原密码', trigger: 'blur' }],
    newPassword: [
        { required: true, min: 6, message: '密码长度不少于6位', trigger: 'blur' },
        { pattern: /^(?=.*[a-zA-Z])(?=.*\d)/, message: '密码需包含字母和数字', trigger: 'blur' }
    ],
    confirmPassword: [
        { required: true, message: '请确认新密码', trigger: 'blur' },
        { validator: (rule, value, callback) => {
            if (value !== form.newPassword) callback(new Error('两次密码不一致'))
            else callback()
        }, trigger: 'blur' }
    ]
}
\end{lstlisting}

% ============================================================
\section{全局交互规范}

\subsection{消息提示规范}
\begin{itemize}
    \item 操作成功：ElMessage.success()，绿色，显示2秒；
    \item 操作失败：ElMessage.error()，红色，显示3秒，附带具体错误原因；
    \item 删除确认：统一使用 ElMessageBox.confirm()。
\end{itemize}

\subsection{加载状态规范}
\begin{itemize}
    \item 页面级加载：el-skeleton骨架屏；
    \item 按钮级加载：el-button的 :loading 属性；
    \item 局部刷新：el-table的 v-loading 指令。
\end{itemize}

\subsection{响应式适配断点}
\begin{itemize}
    \item 桌面端（≥1280px）：论文列表4列，编辑器三栏完整；
    \item 平板端（768-1279px）：论文列表2列，编辑器侧面板默认收起；
    \item 移动端（<768px）：论文列表1列，编辑器仅编辑区 + 抽屉式章节列表。
\end{itemize}
